\documentclass[10pt,a4paper]{amsart}
\usepackage[margin=19mm]{geometry}
\usepackage{amsmath,amssymb,mathtools}
\usepackage[scr=rsfs,cal=euler]{mathalfa}
\usepackage{xfrac}
\usepackage[unicode,hidelinks,bookmarks=false]{hyperref}
\newcommand{\ie}{i.e.\ }
\newcommand{\cf}{cf.\ }
\newcommand{\resp}{respectively\ }
\newcommand{\Subsection}{\subsection}
\providecommand{\nobreakdash}{\nobreak\hbox{-}}
\setlength{\emergencystretch}{2em}
\multlinegap0pt
\usepackage{bm}
\usepackage{bbm}
\usepackage{dsfont}
\usepackage{xspace}
\usepackage{cancel}
\usepackage{mathtools}
\usepackage{amsthm}
\makeatletter
\@ifundefined{theorem}{\newtheorem{theorem}{Theorem}}{}
\@ifundefined{lemma}{\newtheorem{lemma}[theorem]{Lemma}}{}
\makeatother
\title[An algebraic multiscale preconditioner for large sparse SPD ]{An algebraic multiscale preconditioner for large sparse SPD matrices}
\author{Yingjie Zhou, Shubin Fu, Eric Tsz Shun Chung}
\date{Source: arXiv:2606.04864v1 (CC BY 4.0)}
\begin{document}
\maketitle

\noindent\emph{Source and licence.} An algebraic multiscale preconditioner for large sparse SPD matrices. Version arXiv:2606.04864v1 (CC BY 4.0), distributed under the Creative Commons Attribution 4.0 International licence, as recorded in the arXiv licence metadata for this exact version and shown on the abstract page https://arxiv.org/abs/2606.04864v1. Full rights notice: licenses/arxiv-2606.04864-CC-BY-4.0.txt.\par\smallskip

\noindent\textit{Editorial scope.} Counted unit: the Lemma whose source label is \texttt{None} in the article (source0.tex lines 676--684, source label \texttt{None}), delivered together with the article's own complete original proof of it (source0.tex lines 686--729, one \texttt{proof} environment of 44 lines). No mathematics is written, repaired or re-derived: statement and proof are the article's own text line for line. Required context retained uncounted: none. Every definition, display and label of those lines is retained unchanged. Omitted from this package and not redistributed: 1--675, 685--685, 730--1173. Nothing inside a retained excerpt is omitted or altered: every retained body is delivered exactly as the article's lines are written, so no omission receipt of any kind accompanies this package. Citations: the retained excerpts carry no citation command at all, so no bibliography entry of the article is transcribed and no reference is redistributed. Reference adaptation: none was needed and none was made; every reference inside the delivered unit resolves inside it, so no unresolved reference or citation is produced. Printed slips kept literal: no printed word, symbol, hypothesis or number of the counted statement or of its proof was altered. Layout and class: the article's own class is \texttt{elsarticle} with packages including graphicx, multirow, amsthm, mathrsfs, xcolor, textcomp, manyfoot, booktabs, listings, mathtools, amsmath, amssymb, eucal, mismath; the wrapper is the lane's pinned \texttt{amsart} editorial wrapper at 10pt on A4, which restates the theorem-like environments the retained text needs and adds the article's own macro definitions that the retained text needs (none), with extra wrapper packages (bm, bbm, dsfont, xspace, cancel, mathtools). The delivered numbering is the unit's own. Assets: no retained span contains a figure environment, a graphics-inclusion command, a PDF-page inclusion, a table or a TikZ picture; no image file is copied into this package, the package declares no asset dependency and no image file is redistributed with it. Licence: version arXiv:2606.04864v1 of this article is distributed under the Creative Commons Attribution 4.0 International licence, as recorded in the arXiv licence metadata for this exact version and shown on the abstract page https://arxiv.org/abs/2606.04864v1. A full rights notice is delivered at licenses/arxiv-2606.04864-CC-BY-4.0.txt. Characters: every retained excerpt is pure ASCII, so the delivered engine, XeLaTeX with its default Latin Modern text font, reports no missing character.
  \begin{lemma}
    For each $x\in \mathbb{R}^{n}$, there exists vectors
    $x_{1}, x_{2}, \cdots, x_{N}$ such that
    $x = \sum_{i=1}^{N}\mathsf{R}_{i}^{\intercal}x_{i}$ and
    \[
      \sum_{i=1}^{N}\|x_{i}\|_{\widehat{\mathsf{A}}_{i}}^{2}\leq \|x\|_{\mathsf{A}}
      ^{2}.
    \]
  \end{lemma}

  \begin{proof}
    For each subdomain $\Omega_{i}$, let $x_{i}= \mathsf{R}_{i}x$ be the
    restriction of $x$ to $\Omega_{i}$. Then we trivially have $x = \sum_{i=1}^{N}
    \mathsf{R}_{i}^{\intercal}x_{i}$.

    Recall that $\widehat{\mathsf{A}}_{i}$ is constructed by modifying the diagonal
    entries of the local matrix such that its row sums are zero. For symmetric
    matrices with non-positive off-diagonal entries (M-matrices), the associated
    quadratic form can be written as a sum over edges:
    \[
      \|x_{i}\|_{\widehat{\mathsf{A}}_{i}}^{2}= \frac{1}{2}\sum_{\substack{j,k \in \Omega_i \\ j \neq k}}
      (-\mathsf{A}_{jk}) (x_{j}- x_{k})^{2}.
    \]
    Summing over all subdomains gives the total energy associated with all
    interior edges:
    \[
      \sum_{i=1}^{N}\|x_{i}\|_{\widehat{\mathsf{A}}_{i}}^{2}= \sum_{i=1}^{N}\sum_{\substack{j,k \in \Omega_i \\ j < k}}
      (-\mathsf{A}_{jk}) (x_{j}- x_{k})^{2}.
    \]
    On the other hand, the global energy norm $\|x\|_{\mathsf{A}}^{2}$ includes
    contributions from all edges in the graph as well as diagonal terms:
    \[
      \|x\|_{\mathsf{A}}^{2}= \sum_{j=1}^{n}\left(\sum_{k=1}^{n}\mathsf{A}_{jk}\right
      )x_{j}^{2}+ \sum_{j < k}(-\mathsf{A}_{jk})(x_{j}- x_{k})^{2}.
    \]
    Since $\mathsf{A}$ is a stiffness matrix derived from an elliptic problem,
    we have $-\mathsf{A}_{jk}\ge 0$ for $j \neq k$ and the row sums
    $\sum_{k}\mathsf{A}_{jk}\ge 0$ (by weak diagonal dominance). Decomposing the
    edge sum into interior edges (where both nodes are in the same $\Omega_{i}$)
    and cut edges (where nodes belong to different subdomains), we obtain:
    \[
      \sum_{j < k}(-\mathsf{A}_{jk})(x_{j}- x_{k})^{2}= \sum_{i=1}^{N}\sum_{\substack{j,k \in \Omega_i \\ j < k}}
      (-\mathsf{A}_{jk}) (x_{j}- x_{k})^{2}+ \sum_{\text{cut edges}}(-\mathsf{A}_{jk}
      )(x_{j}- x_{k})^{2}.
    \]
    Comparing the expressions, we see that
    $\sum_{i=1}^{N}\|x_{i}\|_{\widehat{\mathsf{A}}_{i}}^{2}$ corresponds exactly
    to the first term of the global energy decomposition. Since the contributions
    from cut edges and row sums are non-negative, it follows immediately that:
    \[
      \sum_{i=1}^{N}\|x_{i}\|_{\widehat{\mathsf{A}}_{i}}^{2}\le \|x\|_{\mathsf{A}}
      ^{2}.
    \]
  \end{proof}

\end{document}
